\chapter{Conclusiones}

\section{Evaluación de Objetivos y Requisitos}

A lo largo del desarrollo del presente proyecto se establecieron una serie de
objetivos y requisitos técnicos orientados a validar la viabilidad de una
Máquina de Cognición-Ejecución (MCE) funcional. 
Estos objetivos han sido en su mayoría cumplidos de forma satisfactoria:

\begin{itemize}
    \item Se ha definido formalmente una arquitectura cognitiva que permite la
    resolución autónoma de problemas a través de un sistema computacional capaz
    de planificar, decidir y ejecutar acciones.
    \item Se ha implementado una MCE real sobre un microprocesador ARM 
    Cortex-A55, en un sistema operativo Arch Linux sin necesidad de 
    GPUs locales y con comunicación con modelos LLM en la nube.
    \item Se ha validado el funcionamiento del sistema en entornos 
    operativos reales, como ordenadores personales, mediante 
    instrucciones en lenguaje natural.
    \item La arquitectura ha sido implementada íntegramente en Python, con módulos
    desacoplados que permiten sustituir, ampliar o adaptar funcionalidades sin
    necesidad de reestructurar el sistema completo.
    \item Se ha incorporado una interfaz funcional por línea de comandos 
    que permite al usuario interactuar con la MCE sin requerimientos 
    gráficos adicionales.
    \item Los agentes desarrollados han demostrado ser capaces de ejecutar
    acciones complejas como abrir aplicaciones, navegar por Internet y controlar
    el teclado y ratón, todo ello a partir de instrucciones en lenguaje natural.
    \item Toda la implementación se encuentra documentada, versionada y
    accesible para futuros desarrollos o réplicas.
\end{itemize}

Estos resultados confirman que los requisitos mínimos para validar la propuesta
han sido alcanzados, consolidando a la MCE como una arquitectura funcional que
puede ser escalada o adaptada en futuras aplicaciones.

\section{Conclusiones Generales}

La implementación de la Máquina de Cognición-Ejecución (MCE) ha demostrado que
es posible replicar una arquitectura cognitiva generalista, capaz de integrar
razonamiento, memoria, planificación y ejecución, en un entorno computacional
limitado y no controlado. Este logro adquiere especial relevancia por su
capacidad de operar sin requerir hardware especializado, ofreciendo un marco
operativo accesible, replicable y científicamente fundamentado.

Más allá de la implementación práctica, uno de los resultados más significativos
de este proyecto es el desarrollo de un fundamento teórico formalizado
con suficiente solidez, rigor y generalidad como para sustentar la publicación
de un artículo científico original, actualmente en proceso de revisión bajo el
título: \textit{"Thinking is not Enough: The B* Expansion Technique for
Enhancing Autonomous LLM Agents"} \citet{bstar-expansion-2025}. Este trabajo, 
derivado directamente de los
principios establecidos en este TFG, propone una capa de optimización centrada
en las acciones (no en el razonamiento) para agentes basados en LLMs,
formalizando matemáticamente el crecimiento de su capacidad resolutiva en
función de un espacio de acciones expandido $B^*$.

Además, se ha validado empíricamente que un enfoque basado en capas cognitivas
desacopladas permite experimentar con distintas estrategias de razonamiento,
percepción o ejecución sin comprometer la estabilidad global del sistema. Esta
propiedad convierte a la MCE en un entorno ideal para experimentación en
cognición computacional, simulación de agentes autónomos y prototipado de
arquitecturas agénticas.

Finalmente, el agente DarkVFR (el agente más eficiente entre los evaluados) ha
evidenciado capacidades avanzadas de adaptación a múltiples entornos operativos,
desde navegadores web hasta aplicaciones de escritorio, logrando resolver tareas
complejas con un consumo de recursos y tokens entre 7 y 12 veces inferior al de
arquitecturas comparables en el estado del arte. Estos resultados refuerzan la
viabilidad de la expansión $B^*$ y la MCE como un marco robusto para el
desarrollo de agentes autónomos.

\section{Competencias Adquiridas}

Antes de comenzar este Trabajo de Fin de Grado, ciertas asignaturas del grado de
Ingeniería Robótica Software fueron clave para cimentar las bases teóricas y
prácticas que permitirían su posterior desarrollo. En particular, la asignatura
de \emph{Diseño de Software} fue fundamental para concebir una arquitectura
modular, robusta y extensible, gracias al uso de patrones de diseño como
mediadores, capas de abstracción y técnicas de programación orientada a objetos
que facilitaron la separación de responsabilidades entre componentes. Por otro lado,
\emph{Sistemas Operativos} aportó los conocimientos necesarios para
interactuar directamente con el sistema operativo: desde la gestión de procesos
y de multithreading hasta la conocimientos teóricos como lazy loading o enlazado de lenguaje C,
lo cual fue esencial tanto para la capa de ejecución como para el diseño del lenguaje
\textit{Exelent}. Finalmente, la asignatura de \emph{Álgebra} permitió 
abordar el desarrollo formal del modelo teórico con rigor matemático, 
facilitando la comprensión de conceptos como espacios, bases y generadores, 
herramientas clave para definir formalmente el espacio de problemas, soluciones y 
acciones, así como la definicion canónica de la MCE.

Durante el desarrollo del proyecto, se han adquirido y reforzado una serie de
competencias técnicas y metodológicas:

\begin{itemize}
    \item \textbf{Diseño de arquitecturas cognitivas:} Formalización, implementación y 
    evaluación de arquitecturas orientadas al razonamiento autónomo. (Con técnicas de
    Prompt Engineering, grafos agénticos, etc.)
    \item \textbf{Programación avanzada en Python:} Uso de
    bibliotecas para IA, visión por computador, gestión de eventos o técnicas
    avanzadas de meta-programación.
    \item \textbf{Despliegue en hardware embebido:} Configuración de sistemas
    operativos ligeros (Arch Linux) en microprocesadores ARM y quemado de bootloader/imagenes.
    \item \textbf{Integración de LLMs en sistemas a tiempo real:} Mediante LangChain, APIs
    externas y servicios en la nube.
    \item \textbf{Desarrollo de interfaces de usuario:} Implementación de interfaces
    basadas en Streamlit y línea de comandos para interactuar con la MCE.
\end{itemize}


\section{Futuras Líneas de Investigación}

Este trabajo ha presentado una arquitectura generalista viable para la
interacción agente-entorno basada en lenguaje natural. Sin embargo, su
desarrollo ha abierto nuevas preguntas que pueden guiar futuras
investigaciones:

\begin{itemize}
    \item \textbf{Optimización del espacio de acciones}: Aunque la MCE ha
    demostrado ser capaz de ejecutar acciones complejas, la generación de
    acciones a partir de instrucciones en lenguaje natural puede ser mejorada
    mediante técnicas de optimización del espacio de acciones. Esto podría incluir
    la creación de nuevas funcionalidades que permitan al agente identificar el
    estado del entorno de una forma más enriquecida o la intreacción simplificada
    con aplicaciones específicas.

    \item \textbf{Aprendizaje basado en experiencia contextualizada:} 
    La arquitectura actual emplea un modelo cognitivo fijo para planificar 
    acciones, pero no aprende activamente a partir de la experiencia. Explorar 
    estrategias que permitan a la MCE construir representaciones internas 
    de éxito/fracaso y utilizarlas para adaptar su planificación futura 
    sería un paso clave hacia el aprendizaje continuo.
\end{itemize}

Estas líneas de investigación no solo amplían el alcance técnico de la MCE, sino
que permitirán avanzar hacia arquitecturas cognitivas verdaderamente
generalistas, capaces de operar de forma robusta en una amplia gama de entornos
dinámicos y abiertos.

